\subsection{CTR: de la ejecución 90A al reemplazo de una clave}
\label{sec:ctr-hybrid-lean}
El módulo \nolinkurl{CTRHybrid.lean} conecta el procedimiento 90A
con el cálculo de probabilidades del módulo \nolinkurl{Reduction.lean}.
Fijamos las monedas del observador y una política que elige el dato
adicional en función de las salidas anteriores que puede observar. El
evaluador de la predicción conserva los bloques ocultos y, para .7,
puede conceder el estado terminal al observador. Las decisiones de éste
sólo utilizan su vista autorizada. La fuerza se expresa como probabilidad
de éxito de un algoritmo con un presupuesto temporal declarado.

\paragraph{1. Operaciones que recibe el oráculo.}
La rama sin dato adicional contiene un tramo de cuatro cifrados.
La rama con dato contiene un tramo de tres seguido del tramo de cuatro.
Para usar el mismo denominador en el cálculo, se puede añadir una cuarta
respuesta ficticia al tramo de tres y descartarla. Los cuatro puntos de
contador son distintos incluso al envolver. Sea \(x=(x_1,x_2,x_3,x_4)\)
la respuesta en esos puntos y \(D\) la máscara ya fijada. Las dos funciones
de continuación reciben, respectivamente,
\[
 \begin{aligned}
   &\bigl(\varnothing,(x_1\|x_2\|x_3)\oplus D\bigr),\\
   &\bigl(x_1,(x_2\|x_3\|x_4)\oplus D\bigr).
 \end{aligned}
\]
\proofat{CTRHybrid.actual_segments}{actual\_\allowbreak{}segments} y \proofat{CTRHybrid.executions_refine}{executions\_\allowbreak{}refine} demuestran
la igualdad de salida y estado con 90A por petición y por ejecución
adaptativa. \proofat{CTRHybrid.at_most_two_segments}{at\_\allowbreak{}most\_\allowbreak{}two\_\allowbreak{}segments} da el máximo de \(2N\)
tramos en \(N\) peticiones. \proofat{CTRHybrid.concrete_request_test}{concrete\_\allowbreak{}request\_\allowbreak{}test} y
\proofat{CTRHybrid.concrete_post_test}{concrete\_\allowbreak{}post\_\allowbreak{}test} comprueban que la continuación real puede
ejecutarse desde las respuestas del oráculo, incluidas las peticiones
posteriores con sus propias claves. El simulador conoce el siguiente
estado numérico. Al oráculo se accede exclusivamente por cifrado bajo su
clave secreta \cite[§§10.2.1.2 y 10.2.1.5.1]{nist90a}.

\paragraph{2. Por qué la próxima clave es uniforme en el prefijo ideal.}
Con cuatro respuestas uniformes independientes, separamos la palabra
descartada o divulgada de los tres bloques que forman el próximo estado.
\proofat{CTRHybrid.independent_words_coordinate_bijection}{independent\_\allowbreak{}words\_\allowbreak{}coordinate\_\allowbreak{}bijection} comprueba las dos
direcciones de esa reagrupación. Para cada historia, palabra separada y
máscara, cada estado final \((k,v)\) admite una única terna inicial:
\((k\|v)\oplus D\).
\proofat{CTRHybrid.ideal_renews_key}{ideal\_\allowbreak{}renews\_\allowbreak{}key} comprueba el conteo uniforme para cada
\((k,v)\). En particular, al fijar \(v\) y lo ya divulgado en el
prefijo ideal, la clave del siguiente tramo sigue siendo uniforme. En el
primer tramo, esa condición debe venir de la siembra, según
\eqref{eq:ctr-seed-closeness} o mediante la transferencia por razón de
densidades del análisis de predicción.

\paragraph{3. Hipótesis criptográfica y consecuencia.}
Sea \(M=2^{128}\), \(B=M^4\), \(D_4=(M)_4=M(M-1)(M-2)(M-3)\).
Para una historia y un contador públicos fijados, el test local puntúa las
continuaciones que permiten acertar el objetivo. Bajo permutación aleatoria
se promedian sus puntuaciones sobre las \(D_4\) respuestas ordenadas
distintas. Bajo función aleatoria se promedian sobre las \(B\) respuestas.
Las puntuaciones son no negativas, luego la suma restringida es menor o
igual que la suma completa. \proofat{CTRHybrid.switching_for_ctr_continuation}{switching\_\allowbreak{}for\_\allowbreak{}ctr\_\allowbreak{}continuation}
verifica esa desigualdad para la continuación del perfil. La hipótesis
PRP de clave uniforme permite sustituir el cifrado real por la permutación
con ventaja a lo sumo \(\alpha_4(T)\). El presupuesto \(T\) incluye
evaluar la continuación y gestionar los datos auxiliares admitidos.


\paragraph{4. Cierre por la ley total, incluidas entradas adaptativas.}
Sea \(S=M|\mathcal K|\) la cardinalidad del estado y
\(\delta=D_4|\mathcal K|\alpha_4(T)\) la holgura local expresada en
unidades de conteo (puede redondearse hacia arriba). Para \(n\) tramos
restantes e historial \(h\), \(R_n(h)\) suma el indicador de éxito de
la ejecución CTR real sobre cada par inicial \((k,v)\).
Para la continuación tras observar una palabra ideal \(y\), escribimos
\(h_y\). En preparación la palabra aislada se descarta y \(h_y\) es el
mismo para los \(M\) valores. El conteo ideal \(I_n(h)\) suma sobre los
\(n\) términos aislados y el estado final uniforme. En particular,
\(R_0(h)=I_0(h)\), e \(I_{n+1}(h)=\sum_y I_n(h_y)\).
El oráculo produce los mismos cuatro bloques del procedimiento real.
Se eliminan los indicadores de respuesta repetida al tomar una cota
superior de éxito; se reindexan las respuestas independientes por
\((y,k',v')\) y se suman los contadores anteriores \(v\). De aquí,
\[
  D_4 R_{n+1}(h)\le M\delta+S\sum_y R_n(h_y).
\]
\proofat{CTRHybrid.reply_reindex}{reply\_reindex} comprueba el cambio de
variables y \proofat{CTRHybrid.derived_step}{derived\_step} deriva la
desigualdad del test PRP local. En la preparación omitida para entrada
ausente, el estado real se conserva y el mismo paso usa únicamente
\(D_4\le B\). Por inducción y normalización se obtiene
\[
 P_{\rm real}\ \le\ \gamma^J P_{\rm ideal}
                 +\alpha_4(T)\sum_{i=0}^{J-1}\gamma^i,
 \qquad \gamma=\frac{M S}{D_4}=\frac{M^4}{(M)_4}.
\]
\proofat{CTRHybrid.direct_epoch}{direct\_epoch} verifica la recurrencia
entera equivalente. \proofat{CTRHybrid.autonomous_end_to_end}{autonomous\_end\_to\_end}
la instancia con \(J=N\) en la rama autónoma.
\proofat{CTRHybrid.mixed_end_to_end}{mixed\_end\_to\_end} usa
\(J=2N\) para una política que elige tanto la presencia como el valor de la
máscara en cada una de las \(N\) llamadas. Un primer tramo omitido consume
cero AES. \proofat{CTRHybrid.concrete_mixed}{concrete\_mixed} verifica que
el experimento real ejecuta exactamente \emph{Generate} en ambos casos.
\proofat{CTRHybrid.ideal_mixed}{ideal\_mixed} obtiene en el experimento ideal
\(N\) bloques uniformes independientes y un estado terminal uniforme
independiente. Estas instancias construyen los pesos de todos los historiales.
Para un objetivo de \(b\) bits ocultos, el conteo ideal da \(P_{\rm ideal}
=2^{-b}\) cuando la información entregada al observador cumple su restricción.
La evidencia del producto debe identificar esa vista, la entropía de
la semilla, la época y el coste de las continuaciones. El kernel comprueba
la reducción cuantitativa desde una hipótesis PRP sobre una clave AES.
